\section{Event graphs}
\label{sec:eventgraphs}

The event graph is the semantic IR produced by core compilation.  It is
small on purpose.  It does not store source names, source occurrences,
rank metadata, custom patch legality predicates, or action kernels
parallel to node semantics.  Those are either source-layer data or
derived from the node row.

\subsection{Graph data}
\label{sec:eg:data}

An event graph is a pair of lists:
\[
  \EG = (\mathsf{initialFields},\mathsf{nodes}).
\]
An initial field row contains a type, an optional owner, and an initial
value.  An event node row contains the type and owner of its output
field together with a semantic payload:
\[
  \mathsf{NodeSem} ::= \mathsf{sample}(D)
    \mid \mathsf{commit}(p,R)
    \mid \mathsf{reveal}(f).
\]
Initial field ids are \(0,\ldots,|\mathsf{initialFields}|-1\).  The
output field of node \(n\) is
\[
  \mathsf{nodeTarget}(n) = |\mathsf{initialFields}| + n.
\]
Thus event-field writers are not stored as a second map: the writer of
an event field is determined by arithmetic on the graph row.

The node payloads use typed field references.  Payload evaluation is not
given the whole store.  Instead it receives a \(\mathsf{ReadEnv}\) for
its declared typed read footprint.  This is the boundary that makes the
read set semantic: graph-local code can only read fields that appear in
its footprint.

\subsection{Reads and well-formedness}
\label{sec:eg:wf}

The read footprint is the scheduling and observation footprint:
\begin{itemize}\itemsep1pt
\item a sample reads the public fields declared by its distribution;
\item a commit choice footprint contains the fields visible to the
  committing player at the commit site;
\item a reveal reads the hidden source field being opened.
\end{itemize}
For commits this footprint is deliberately the information available to the
choice, not the minimal syntactic dependency set of the guard expression.  A
guard may ignore some of that environment; scheduling and player observation
still use the larger choice footprint, because strategies may condition on the
whole visible commit-site view.
Graph well-formedness is a proposition.  For each node, every read field
must already be available: it is either initial or written by an earlier
node.  Public internal computations may read only public fields.  Commit
choice footprints must be public or owned by the committing player.  A reveal must
open a hidden field into a public output field of the same type.

Graph reads are not raw SSA ancestry.  They are the information dependencies
that can affect scheduling and player observation.  Pure administrative source
lets are erased before this layer, so they do not create observable graph
nodes.

\subsection{Canonical prerequisites}
\label{sec:eg:prereqs}

Prerequisites are not stored in graph data.  They are derived
canonically:
\[
  q \in \prereq(n)
  \quad\Longleftrightarrow\quad
  q < n \;\wedge\;
  \bigl(\mathsf{nodeTarget}(q) \in \readsOf(n)
    \;\vee\;
    n\text{ is a reveal and }q\text{ is a commit}\bigr).
\]
The first disjunct is the ordinary information dependency: an event
that reads a field waits for the event that wrote the field.  The
second disjunct is the commit/reveal barrier.  It makes a reveal wait
for prior commits even when the reveal does not read their payloads.
This is what enforces the simultaneous-commit discipline in protocols
such as sealed bidding and Matching Pennies.

The negative cases are just as important.  Commit/commit pairs are
independent unless the later commit choice footprint contains a field written
by the earlier one.  Reveal/reveal pairs are independent at this abstraction
level unless one reveal reads a field written by the other, which the
current core does not generate.  Public samples do not wait for hidden
commitments.

\subsection{Configurations and availability}
\label{sec:eg:configurations}

A configuration stores a finite set of completed nodes and a typed field
store:
\[
  C = (\mathsf{done},\mathsf{store}).
\]
The initial configuration has no completed nodes and contains exactly the
initial field values.  Completing a node records the node as done and
writes its canonical output field.  There is no operation that completes
a node while writing an arbitrary field.

A node is ready when it is not done and all canonical prerequisites are
done.  A commit action is available when the node is ready, owned by the
acting player, the action value has the guard type, the guard's choice
environment can be built from the current store, and the guard evaluates
to true.  An internal event is available when a ready sample or reveal can
evaluate from the current store.

\subsection{Schedule independence}
\label{sec:eg:diamond}

The graph layer proves local schedule-independence facts.  If two
distinct available events are ready at the same configuration, completing
one preserves availability of the other, because ready nodes do not read
each other's output fields.  Supported results for the two events can be
swapped to reach the same raw graph configuration.

These theorems do not assert that every real runtime hides every
implementation trace.  They say that the graph state and graph
observations are functions of the completed-node downset and field
values, not of an irrelevant ordering witness.  Strategic presentations
can therefore quotient or checkpoint primitive order without changing
the graph-level state they observe.

\begin{theorem}[Frontier diamond]
\label{thm:diamond}
For two distinct supported available primitive events at the same
reachable configuration, there are two two-event batches, one executing
the events in each order, whose final raw graph configurations are
equal.  Consequently the public and player graph observations at the two
checkpoints are equal.
\end{theorem}

\begin{theorem}[Commit/reveal barrier]
\label{thm:barrier}
If a reveal node is ready, every earlier commit node is already complete.
Equivalently, a prior commit action is not available at a checkpoint where
the reveal is ready.
\end{theorem}

\subsection{Compiler adequacy at the graph boundary}
\label{sec:eg:adequacy}

Compilation produces a \(\mathsf{CompiledProgram}\) containing:
\begin{itemize}\itemsep1pt
\item a well-formed graph;
\item payoff projections as graph-local \(\mathsf{EventPayoff}\)
  entries;
\item a proof that payoff reads are public and available by graph
  termination.
\end{itemize}
For terminal reachable primitive-machine states, payoff projection is
total and agrees with the source payoff expressions reconstructed from
the terminal compiled store.

\paragraph{Lean correspondence.}
Graph data and canonical prerequisites are in
\TT{Vegas.EventGraph.Basic}.  Primitive execution, readiness,
observations, and schedule-diamond lemmas are in
\TT{Vegas.EventGraph.Execution} and \TT{Vegas.EventGraph.Batch}.  The
compiler is \TT{Vegas.Core.ToEventGraph.Compiler}; source payoff
adequacy is \TT{Vegas.Core.ToEventGraph.SourceAdequacy} and
\TT{Vegas.Core.Theorems.Outcome}.
